\begin{center}
\begin{tikzpicture}[
    thick,
    % Stile per le linee dei fermioni (linea con freccia a metà)
    fermion/.style={draw=black, postaction={decorate},
        decoration={markings, mark=at position 0.55 with {\arrow{latex}}}},
    % Stile per i bosoni W (linea ondulata)
    boson/.style={draw=black, decorate,
        decoration={snake, amplitude=1mm, segment length=3.5mm}},
    % Stile per i vertici dei nucleoni (i cerchi grigi)
    blob/.style={circle, draw=black, fill=gray!30, minimum size=6mm, inner sep=0pt}
]

% ---------------- NUCLEONE SUPERIORE ----------------
% Frecce esterne n e p
\draw[fermion] (-4.5, 2) node[left, font=\Large] {n} -- (-3, 2);
\draw[fermion] (3, 2) -- (4.5, 2) node[right, font=\Large] {p};

% Le 3 linee dei quark all'interno (tra i due blob)
\draw[fermion] (-3, 2.15) -- (3, 2.15);
\draw[fermion] (-3, 2.0) -- (3, 2.0);
\draw[fermion] (-3, 1.85) -- (3, 1.85);

% Bosone W- 
\draw[boson] (0, 1.85) -- (1.5, 0.8);
\node[font=\Large] at (0.4, 1.2) {$W^-$};

% Elettrone e antineutrino 
% (Nota: la linea del neutrino è disegnata da destra a sinistra per invertire la freccia)
\draw[fermion] (1.5, 0.8) -- (3.5, 1.3) node[right, font=\Large] {$e^-$};
\draw[fermion] (3.5, 0.3) node[right, font=\Large] {$\overline{\nu}_e$} -- (1.5, 0.8);


% ---------------- NUCLEONE INFERIORE ----------------
% Frecce esterne n e p
\draw[fermion] (-4.5, -2) node[left, font=\Large] {n} -- (-3, -2);
\draw[fermion] (3, -2) -- (4.5, -2) node[right, font=\Large] {p};

% Le 3 linee dei quark all'interno (tra i due blob)
\draw[fermion] (-3, -1.85) -- (3, -1.85);
\draw[fermion] (-3, -2.0) -- (3, -2.0);
\draw[fermion] (-3, -2.15) -- (3, -2.15);

% Bosone W- 
\draw[boson] (0, -1.85) -- (1.5, -0.8);
\node[font=\Large] at (0.4, -1.2) {$W^-$};

% Elettrone e antineutrino
\draw[fermion] (1.5, -0.8) -- (3.5, -0.3) node[right, font=\Large] {$e^-$};
\draw[fermion] (3.5, -1.3) node[right, font=\Large] {$\overline{\nu}_e$} -- (1.5, -0.8);


% --- DISEGNO I BLOB ---
% (Disegnati per ultimi così coprono perfettamente le estremità delle linee)
\node[blob] at (-3, 2) {};
\node[blob] at (3, 2) {};
\node[blob] at (-3, -2) {};
\node[blob] at (3, -2) {};

\end{tikzpicture}
\end{center}